% GENERATED FILE. Edit canonical lessons, then run npm run generate:books.
% canonical-source-hash: fnv1a64:78296253fe250023
% canonical-lessons: FR-C393-garantie, FR-C393-garantir, FR-C393-devis, FR-C393-rentable, FR-C393-avantageux

\chapter{Garanties et devis}
\label{ch:fr-garanties-et-devis}
\hlchaptermodality{393}

\begin{chapteropening}
\textbf{By the end of this chapter:} \emph{I can deal with what I buy: a guarantee, guaranteeing something, a quote, and whether it pays or is good value.}

Complete the last lesson of chapter 393: I can deal with what I buy: a guarantee, guaranteeing something, a quote, and whether it pays or is good value.
\end{chapteropening}

\section[la garantie]{la garantie — the guarantee, the warranty}

\label{lesson:FR-C393-garantie}



\begin{quote}
Before the new one: say the French for the bank transfer, then the French for to take out.
\end{quote}

\subsection*{You'll want to know: la garantie}
\textbf{la garantie} — \textquotedblleft{}the guarantee, the warranty\textquotedblright{}.

\begin{quote}
Le frigo est encore sous garantie.
\end{quote}

\begin{quote}
\emph{The fridge is still under warranty.}
\end{quote}

\emph{Sous garantie} — under warranty.

\begin{cousinweb}
\emph{garantie} is formed on \emph{garantir} (EtymDB), from Old French \emph{garant} \textquotedblleft{}protector, warrant\textquotedblright{}, usually traced to a Frankish word.
\end{cousinweb}

\begin{grammarlens}[title={What you've built}]
One more word for telling, explaining and describing at B1.
\end{grammarlens}

\subsection*{Your turn}
\begin{itemize}
\raggedright
  \item \emph{Say it:} \emph{la garantie}
  \item \emph{Say it:} \emph{la garantie}, once more
  \item \emph{Say it:} a sentence of your own with \emph{la garantie}
  \item \emph{Recall:} say the French for the bank transfer, then the French for to take out, then say \emph{la garantie} again
\end{itemize}

\begin{tcolorbox}[breakable,colback=teal!4,colframe=teal!35!black,title={Before you move on}]
What does la garantie mean? (\textquotedblleft{}The guarantee, the warranty\textquotedblright{}.) Say it once more.
\end{tcolorbox}
\section[garantir]{garantir — to guarantee}

\label{lesson:FR-C393-garantir}



\begin{quote}
Before the new one: say the French for to take out, then the French for the guarantee.
\end{quote}

\subsection*{You'll want to know: garantir}
\textbf{garantir} — \textquotedblleft{}to guarantee\textquotedblright{}.

\begin{quote}
Le vendeur garantit la réparation pendant un an.
\end{quote}

\begin{quote}
\emph{The seller guarantees the repair for a year.}
\end{quote}

A \emph{-ir} verb like \emph{finir}.

\begin{cousinweb}
\emph{garantir} is inherited from Old French \emph{garantir} \textquotedblleft{}to protect, to vouch for\textquotedblright{} (EtymDB), from \emph{garant} \textquotedblleft{}protector\textquotedblright{}, usually traced to a Frankish word.
\end{cousinweb}

\begin{grammarlens}[title={What you've built}]
One more word for telling, explaining and describing at B1.
\end{grammarlens}

\subsection*{Your turn}
\begin{itemize}
\raggedright
  \item \emph{Say it:} \emph{garantir}
  \item \emph{Say it:} \emph{garantir}, once more
  \item \emph{Say it:} a sentence of your own with \emph{garantir}
  \item \emph{Recall:} say the French for to take out, then the French for the guarantee, then say \emph{garantir} again
\end{itemize}

\begin{tcolorbox}[breakable,colback=teal!4,colframe=teal!35!black,title={Before you move on}]
What does garantir mean? (\textquotedblleft{}To guarantee\textquotedblright{}.) Say it once more.
\end{tcolorbox}
\section[le devis]{le devis — the quote, the estimate}

\label{lesson:FR-C393-devis}



\begin{quote}
Before the new one: say the French for the guarantee, then the French for to guarantee.
\end{quote}

\subsection*{You'll want to know: le devis}
\textbf{le devis} — \textquotedblleft{}the quote, the estimate\textquotedblright{}.

\begin{quote}
Le plombier m'a envoyé un devis.
\end{quote}

\begin{quote}
\emph{The plumber sent me a quote.}
\end{quote}

The \emph{s} is silent. \emph{Demander un devis} — to ask for a quote.

\begin{cousinweb}
\emph{devis} is formed on \emph{deviser} \textquotedblleft{}to divide, to arrange\textquotedblright{} (EtymDB), from a Vulgar Latin form of \emph{dīvidere} \textquotedblleft{}to divide\textquotedblright{} (Lewis \& Short).
\end{cousinweb}

\begin{grammarlens}[title={What you've built}]
One more word for telling, explaining and describing at B1.
\end{grammarlens}

\subsection*{Your turn}
\begin{itemize}
\raggedright
  \item \emph{Say it:} \emph{le devis}
  \item \emph{Say it:} \emph{le devis}, once more
  \item \emph{Say it:} a sentence of your own with \emph{le devis}
  \item \emph{Recall:} say the French for the guarantee, then the French for to guarantee, then say \emph{le devis} again
\end{itemize}

\begin{tcolorbox}[breakable,colback=teal!4,colframe=teal!35!black,title={Before you move on}]
What does le devis mean? (\textquotedblleft{}The quote, the estimate\textquotedblright{}.) Say it once more.
\end{tcolorbox}
\section[rentable]{rentable — profitable, worth it}

\label{lesson:FR-C393-rentable}



\begin{quote}
Before the new one: say the French for to guarantee, then the French for the quote.
\end{quote}

\subsection*{You'll want to know: rentable}
\textbf{rentable} — \textquotedblleft{}profitable, worth it\textquotedblright{}.

\begin{quote}
Ce magasin n'est plus rentable.
\end{quote}

\begin{quote}
\emph{This shop is no longer profitable.}
\end{quote}

From \emph{rente}, income.

\begin{cousinweb}
\emph{rentable} is built from \emph{rente} \textquotedblleft{}income\textquotedblright{} and \emph{-able}. \emph{rente} comes from a Vulgar Latin \emph{rendita}, from the verb behind \emph{rendre} \textquotedblleft{}to give back\textquotedblright{}, Latin \emph{reddere} (Lewis \& Short).
\end{cousinweb}

\begin{grammarlens}[title={What you've built}]
One more word for telling, explaining and describing at B1.
\end{grammarlens}

\subsection*{Your turn}
\begin{itemize}
\raggedright
  \item \emph{Say it:} \emph{rentable}
  \item \emph{Say it:} \emph{rentable}, once more
  \item \emph{Say it:} a sentence of your own with \emph{rentable}
  \item \emph{Recall:} say the French for to guarantee, then the French for the quote, then say \emph{rentable} again
\end{itemize}

\begin{tcolorbox}[breakable,colback=teal!4,colframe=teal!35!black,title={Before you move on}]
What does rentable mean? (\textquotedblleft{}Profitable, worth it\textquotedblright{}.) Say it once more.
\end{tcolorbox}
\section[avantageux]{avantageux — good value, favourable}

\label{lesson:FR-C393-avantageux}



\begin{quote}
Before the new one: say the French for the quote, then the French for profitable.
\end{quote}

\subsection*{You'll want to know: avantageux}
\textbf{avantageux} — \textquotedblleft{}good value, favourable\textquotedblright{}.

\begin{quote}
Cet abonnement est très avantageux.
\end{quote}

\begin{quote}
\emph{This subscription is very good value.}
\end{quote}

Feminine \emph{avantageuse}.

\begin{cousinweb}
\emph{avantageux} is inherited from Old French \emph{avantageus} (EtymDB), built on \emph{avantage} \textquotedblleft{}advantage\textquotedblright{}, from \emph{avant} \textquotedblleft{}ahead\textquotedblright{}.
\end{cousinweb}

\begin{grammarlens}[title={What you've built}]
One more word for telling, explaining and describing at B1.
\end{grammarlens}

\subsection*{Your turn}
\begin{itemize}
\raggedright
  \item \emph{Say it:} \emph{avantageux}
  \item \emph{Say it:} \emph{avantageux}, once more
  \item \emph{Say it:} a sentence of your own with \emph{avantageux}
  \item \emph{Recall:} say the French for the quote, then the French for profitable, then say \emph{avantageux} again
\end{itemize}

\begin{tcolorbox}[breakable,colback=teal!4,colframe=teal!35!black,title={Before you move on}]
What does avantageux mean? (\textquotedblleft{}Good value, favourable\textquotedblright{}.) Say it once more.
\end{tcolorbox}
